\documentclass{article}

\usepackage[a4paper]{geometry}
\usepackage[utf8]{inputenc}
\usepackage[english]{babel}
\usepackage[noend]{algpseudocode}
\usepackage{amsthm}
\usepackage{amsmath}
\usepackage{colortbl}
\usepackage{graphicx}
\usepackage{lineno}
\usepackage[textwidth=4cm]{todonotes}
\usepackage{changepage}
\usetikzlibrary{tikzmark}
\usepackage{hyperref}

\usepackage[backend=biber]{biblatex}
\usepackage{tcolorbox}
\tcbuselibrary{breakable}

\addbibresource{references.bib}

\geometry{left=3cm, right=6cm, top=3cm, bottom=3cm, marginparwidth=4cm} 

\setuptodonotes{noline, color=green!15, size=\scriptsize}

\newtheorem*{exercise}{Exercise}
\newcommand{\grey}{\cellcolor{black!25}}

\title{Algorithms and Data Structures \\
  Handin 1: Searching}

\author{
    202605673 Lukas Alstrup \\
    202605116 Thomas Seeberg Hansen \\
    202606296 Frederik Dommer Løvbjerg Plum
}

\date{August 31, 2026}

\begin{document}

\maketitle
\begin{exercise}[Searching]
	Assume we have an infinite long increasing sequence of real values $x_1 < x_2 < x_3 < ...$ and we want to find the position of a real value
	$y > x_1$ in this list. More specifically, we want to find the index $d$ where $x_{d - 1} < y \leq x_d$, i.e.
	$x_d = y$ if $y$ is in the list, and otherwise $x_d$ is the successor of $y$ in the list. E.g. for the sequence $3,
		5, 7, 9, 11, 17, 19, 23, 31, 33, 37, ...$ and $y = 11$ we have $d = 5$, since $y = x_5 = 11$, whereas for $y = 24$
	we have $d = 9$, since $23 = x_8 \leq y \leq x_9 = 31$. We assume the sequence to be divergent, i.e. there
	always exists such a $d$.
	Describe an algorithm to find the index $d$, where the number of comparisons performed
	between $y$ and the $x_is$ is a function of the final value of d. What is the most efficient algorithm
	(fewest number of comparisons performed) you can find? Argue that your algorithm finds the
	correct index d, such that either $y = x_d$ or $x_{d - 1} < y < x_d$. Analyze your algorithm — how many
	comparisons does your algorithm perform as a function of $d$?
	(This problem is intentionally open ended; it was studied by Bentley and Yao in 1976 \cite{BENTLEY197682} if
	you are interested in the research on the problem).
\end{exercise}

\section*{Comment}
\label{sec:comment}
To find the index $d$ of a real value $y$ in an infinite long increasing sequence of real values $A=\{x_1 < x_2<x_3...\}$, a two-phased approach with both exponential search and binary search is performed due to an unknown upper bound for $A$. For the created algorithm it is assumed that the maximum value of $A$ cannot be accessed, otherwise a single-phased binary search algorithm would be sufficient to find $d$ based on $y$.

\section*{Solution}
A simple (and quite effective) solution to deal with the unknown upper bound caused by the infinite length of $A$ (see \nameref{sec:comment}) is to split the algorithm into a two-phased approach:
\begin{enumerate}
	\item Determining a subset $B$ such that $B \subseteq A$ using an exponential "doubling" technique.
	\item Performing binary search on $B$.
\end{enumerate}

\subsection*{Phase 1. Exponential search / doubling}
In this phase we find the subset of A in which y is located. This is done by doubling a search index $i$ until value of element at index $i$ is greater or equal to our desired value $y$.




%Rewrite this
% During the first phase we search $A$ in order to be able to determine a lower and upper boundary for the binary search. For this purpose an exponential search is used which iterates over the sequence $A$ by incrementing the index by double what it was in the previous iteration, hence making the range being searched exponentially larger on each iteration.
% We can express this by a variable $k$ which is used to track the amount of doublings that has been performed so far since the probe indices are a power of 2, $2^k$ for $k \in \{0, 1, 2, 3, ...\}$.
% The exponential search must continue until $y \leq 2^k$ which means that at that point $y > 2^{k-1}$.

% \begin{equation}
% 	2^{k-1} < d \leq 2^{k}
% \end{equation}

% We can then take the binary logarithm
% \begin{equation}
% 	\log_2{2^{k-1}} < \log_2{d} \leq \log_2{2^k} \iff k-1 < \log_2{d} \leq k \iff \lceil \log_2{d} \rceil = k
% \end{equation}

\begin{center}
	%\includegraphics[width=8cm]{whiteboard.jpg}
	\begin{tabular}{ccccccccccccc}
		                                       & 1 & 2 & 3 & 4 & 5            & 6 & 7 & 8 & 9 & 10 & 11 & $\infty$ \\
		\cline{2-13}
		$A$                                    &
		\multicolumn{1}{|c|}{\tikzmark{aa}3}   &
		\multicolumn{1}{|c|}{\tikzmark{bb}5}   &
		\multicolumn{1}{|c|}{\tikzmark{cc}7}   &
		\multicolumn{1}{|c|}{\tikzmark{dd}9}   &
		\multicolumn{1}{|c|}{\tikzmark{ee}11}  &
		\multicolumn{1}{|c|}{\tikzmark{ff}17}  &
		\multicolumn{1}{|c|}{\tikzmark{gg}19}  &
		\multicolumn{1}{|c|}{\tikzmark{hh}23}  &
		\multicolumn{1}{|c|}{\tikzmark{ii}931} &
		\multicolumn{1}{|c|}{\tikzmark{jj}33}  &
		\multicolumn{1}{|c|}{\tikzmark{kk}37}  &
		\multicolumn{1}{|c|}{\tikzmark{ll}...}                                                                     \\
		\cline{2-13}

		                                       &   &   &   &   & $\downarrow$                                      \\
		                                       &   &   &   &   & y
	\end{tabular}
	\begin{tikzpicture}[ remember picture, overlay]
		\draw [->,red] ([yshift=5ex]pic cs:aa) [bend left] to ([yshift=5ex]pic cs:bb);
		\draw [->,red] ([yshift=5ex]pic cs:bb) [bend left] to ([yshift=5ex]pic cs:dd);
		\draw [->,red] ([yshift=5ex]pic cs:dd) [bend left] to ([yshift=5ex]pic cs:hh);

	\end{tikzpicture}
\end{center}
% The mathematical invariant of our exponential search algorithm is
% \begin{equation}
% 	1 \leq i < i\cdot2 \wedge x_i < y \ \text{for all} \ 1 \leq i < d \wedge y < x_i \ \text{for all} \ d < i
% \end{equation}

The pseudo code for our exponential search can be written as
\begin{center}
	\begin{minipage}{8.8cm}
		\begin{algorithmic}
			\Procedure{ExponentialSearch}{$A[x_1,x_2,x_3,..., x_{\infty}], x$}

			\State{$i = 1$}
			\While{$A[i] < y$}
			\State{$i = i\cdot2$}
			\EndWhile
			\State{\Return{i}}
			\EndProcedure
		\end{algorithmic}
	\end{minipage}
\end{center}

This algorithm works because our list $A$ is a sorted list, so when a element of the list is less than an element at index $i$, then our element must be present before $i$, if our element is contained in $A$.
We have therefore with our algorithm found a subset of $A$ which is of finite length and we can therefore use a search algorithm such as binary search, which is much more performant that linear search.

\subsection*{Phase 2: Binary search}
In phase 1 we found a subset $B$ whose largest element is greater or equal to $y$. We therefore know that if the element $y$ is contained in the infinite list $A$ then $B$ must contain it (because $A$ is a sorted list). We can therefore perform a binary search on $B$ and find $y$ if it is contained.
% In the exponential search the value $k$ as a function of of $d$ is determined, and this marks the lower and upper bound that needs to be searched in the binary search.

\begin{center}
	\begin{minipage}{8.8cm}
		\begin{algorithmic}
			\Procedure{BinarySearch}{$B[x_1,x_2,x_3,..., x_n], x$}

			\State{$low = 1$}
			\State{$high = |B|+1$}
			\While{$low < high$}
			\State{$mid = \lfloor(high + low)/2\rfloor$}
			\If{$x=A[mid]$}
			\State{\Return{$mid$}}
			\ElsIf{$x<A[mid]$}
			\State{$high = mid$}
			\ElsIf{$x>A[mid]$}
			\State{$low = mid$}
			\EndIf
			\EndWhile
			\State{\Return{$-1$}}

			\EndProcedure
		\end{algorithmic}
	\end{minipage}
\end{center}

And our visual invariant for binary search would be

\begin{center}
	%\includegraphics[width=8cm]{whiteboard.jpg}
	\begin{tabular}{ccccccccccccc}
		                                                      & 1 & 2 & 3 & 4 & 5            & 6 & 7 & 8 & 9 & 10 & 11 & $n$ \\
		\cline{2-13}
		$B$                                                   &
		\multicolumn{1}{|c|}{\cellcolor{gray} \tikzmark{a}3}  &
		\multicolumn{1}{|c|}{\cellcolor{gray} \tikzmark{b}5}  &
		\multicolumn{1}{|c|}{\tikzmark{c}7}                   &
		\multicolumn{1}{|c|}{\tikzmark{d}9}                   &
		\multicolumn{1}{|c|}{\tikzmark{e}11}                  &
		\multicolumn{1}{|c|}{\tikzmark{f}17}                  &
		\multicolumn{1}{|c|}{\tikzmark{g}19}                  &
		\multicolumn{1}{|c|}{\tikzmark{h}23}                  &
		\multicolumn{1}{|c|}{\cellcolor{gray} \tikzmark{i}31} &
		\multicolumn{1}{|c|}{\cellcolor{gray} \tikzmark{j}33} &
		\multicolumn{1}{|c|}{\cellcolor{gray} \tikzmark{k}37} &
		\multicolumn{1}{|c|}{\cellcolor{gray} \tikzmark{l}...}                                                               \\
		\cline{2-13}

		                                                      &   &   &   &   & $\downarrow$                                 \\
		                                                      &   &   &   &   & y
	\end{tabular}

\end{center}
Our binary search algorithm works because we give it a finite sorted list $B$. When searching we check whether our wanted number $y$ is greater or lower than our middle point. This way we work towards our wanted number $y$ if it exsists in our list $B$ and our original list $A$. If we didn't use a list with a finite amount of elements (like trying to do binary search on our original list $A$) then it would simply not be possible, because we would not have a index for $high$ to start with.


\newpage
\subsection*{The combined algorithm}
Our combined pseudo code would be written as
\begin{center}
	\begin{minipage}{8.8cm}
		\begin{algorithmic}
			\Procedure{Search}{$A[x_1,x_2,x_3,..., x_{\infty}], x$}

			\Comment{Phase 1. Exponential search}

			\State{$low, high = 0, 1$}
			\While{$high < y$}
			\State{$low = high$}
			\State{$high = high \cdot 2$}
			\EndWhile

			\Comment{Phase 2. Binary search}
			\State{$high = high + 1$}
			\While{$low < high$}
			\State{$mid = \lfloor(high + low)/2\rfloor$}
			\If{$x=A[mid]$}
			\State{\Return{$mid$}}
			\ElsIf{$x<A[mid]$}
			\State{$high = mid$}
			\ElsIf{$x>A[mid]$}
			\State{$low = mid$}
			\EndIf
			\EndWhile
			\State{\Return{$-1$}}


			\EndProcedure
		\end{algorithmic}
	\end{minipage}
\end{center}

It follows that the efficiency of the algorithm is dependent on the index $d$ of our wanted element $y$. The two phases of our algorithm have different efficiencies. The index gotten from phase one is $2\cdot1\rightarrow2\cdot2\cdot2\rightarrow2\cdot2\cdot2$ which can be simplified to $2^k$ where $k$ is the number of iterations taken. Our algorithm checks for $d\leq2^k$ and we can therefore solve for $k$ using a logarithm $\log_2{d}\leq k$. The number of iterations given by index $d$ is therefore $\lceil \log_2{d}\rceil=k$.

The second phase of our search is a binary search between the last index $2^{k-1}$ and the current index from phase 1 $2^{k}$. The interval size if therefore $2^{k} - 2^{k-1}$. This can be simplified to $2^{k-1}$. From \textcite{SupplementaryADS} \textbf{Theorem 6} we know that the worst case optimality of binary search can be calculated from $\lceil\log_2{(n+1)}\rceil$ where $n = 2^{k-1}$. We know from phase one that $k = \lceil \log_2{d}\rceil$, and therefore $n = 2^{\lceil \log_2{d}\rceil-1}$

We can now combine our two equations for iterations into one equation of how many iterations are run as a worst case as a function of $d$.
$$f(d)=\lceil \log_2{(d)} \rceil + \lceil\log_2{(2^{\lceil \log_2{d}\rceil-1}+1)}\rceil$$

\noindent\begin{tcolorbox}[breakable, toprule at break=0pt,bottomrule at break=0pt, colback={blue!5!white}, colframe={black!70!white}, arc={0mm}]
	To verify that our function for calculating how many iterations are is correct we'll go through a simple example of the algorithm.

	In this example we'll use the table
	\begin{center}
		%\includegraphics[width=8cm]{whiteboard.jpg}
		\begin{tabular}{ccccccccccccc}
			                                     & 1 & 2 & 3 & 4 & 5 & 6 & 7 & 8 & 9 & 10 & 11 & \\
			\cline{2-12}                         &
			\multicolumn{1}{|c|}{\tikzmark{a}3}  &
			\multicolumn{1}{|c|}{\tikzmark{b}5}  &
			\multicolumn{1}{|c|}{\tikzmark{c}7}  &
			\multicolumn{1}{|c|}{\tikzmark{d}9}  &
			\multicolumn{1}{|c|}{\tikzmark{e}11} &
			\multicolumn{1}{|c|}{\tikzmark{f}17} &
			\multicolumn{1}{|c|}{\tikzmark{g}19} &
			\multicolumn{1}{|c|}{\tikzmark{h}23} &
			\multicolumn{1}{|c|}{\tikzmark{i}31} &
			\multicolumn{1}{|c|}{\tikzmark{j}33} &
			\multicolumn{1}{|c|}{\tikzmark{k}37}                                                 \\
			\cline{2-12}
		\end{tabular}
	\end{center}
	and try to find the element $a = 19$

	We'll go through phase 1 first. In this phase we'll after every step check if the given element $y$ at the index $d$ is $\geq$ to the wanted element $a$. The phase 1 would then look like this:
	$$1 \rightarrow 2 \rightarrow 4 \rightarrow 8$$
	At $d = 8$ we see that $A[8] \geq 19$ and we are therefore finished with phase 1. We have so far run 4 iterations. We will now run phase 2 of the algorithm which is binary search. We'll perform the binary search between index 4 and index 8.

	In this binary search we start with the middle as index $6$. The element at index $6$ is less than our wanted element $a$ and index $6$ becomes our new low. The new middle is now index $7$. Index $7$'s value is $19$ which was exactly the value we're looking for.

	We've now performed $4$ iterations in phase 1 and $2$ iterations in phase 2. The total amount of iterations is therefore $6$.

	We can now  try to see what our function $f(d)$ returns for finding an element at index position $7$.
	$$f(7)=\lceil \log_2{(7)} \rceil + \lceil\log_2{(2^{\lceil \log_2{7}\rceil-1}+1)}\rceil = 6$$
	So our pratical work through of the problem and our equation match up at $6$ iterations!

\end{tcolorbox}

\printbibliography

\end{document}
